%\iffalse
% imac.dtx generated using makedtx version 1.3 (c) Nicola Talbot
% Command line args:
%   -dir imac-source
%   -src imac.sty=>imac.sty
%   -license lppl
%   -doc imac-source/imac.tex
%   -author
%   "Joseph C. Slater"
%   imac
% Created on 2026/10/7 14:46
%\fi
%\iffalse
%<*package>
%% \CharacterTable
%%  {Upper-case    \A\B\C\D\E\F\G\H\I\J\K\L\M\N\O\P\Q\R\S\T\U\V\W\X\Y\Z
%%   Lower-case    \a\b\c\d\e\f\g\h\i\j\k\l\m\n\o\p\q\r\s\t\u\v\w\x\y\z
%%   Digits        \0\1\2\3\4\5\6\7\8\9
%%   Exclamation   \!     Double quote  \"     Hash (number) \#
%%   Dollar        \$     Percent       \%     Ampersand     \&
%%   Acute accent  \'     Left paren    \(     Right paren   \)
%%   Asterisk      \*     Plus          \+     Comma         \,
%%   Minus         \-     Point         \.     Solidus       \/
%%   Colon         \:     Semicolon     \;     Less than     \<
%%   Equals        \=     Greater than  \>     Question mark \?
%%   Commercial at \@     Left bracket  \[     Backslash     \\
%%   Right bracket \]     Circumflex    \^     Underscore    \_
%%   Grave accent  \`     Left brace    \{     Vertical bar  \|
%%   Right brace   \}     Tilde         \~}
%</package>
%\fi
% \iffalse
% Doc-Source file to use with LaTeX2e
% Copyright (C) 2026 Joseph C. Slater, all rights reserved.
% \fi
% \iffalse
%<*driver>
%%%%%%%%%%%%%%%%%%%%%%%%%%% ASME.tex %%%%%%%%%%%%%%%%%%%%%%%%%%%%%%%%%
%%% This is a template for the IMAC (International Modal Analysis
%%% Conference) style file imac.sty.  No subsubsection format is
%%% defined for IMAC, so I used this command for the nomenclature command.
%%% With this you should also find imac.bst which applies the closest
%%% thing I can determine to be the ``Standard'' for IMAC.
%%% The 9 point text size has already been defined in the style file imac.sty.
%%%
%%% Version 2026-10-07
%%% (versioned by release date rather than a version number; see the
%%% "Changes" section below for the full history)
%%% As of the 2026-10-07 release, this style matches the official IMAC
%%% Word Template.dotx: single-column A4 layout, Times New Roman body
%%% text, sentence-case bold headings, and plain (non-superscript)
%%% bracketed citations.
%%%
\documentclass{article}
\usepackage{imac}

\usepackage{hyperref}
\hypersetup{
    colorlinks=true,
    linkcolor=blue,
    citecolor=blue,
    filecolor=magenta,
    urlcolor=cyan,
    pdftitle={IMAC},
    pdfpagemode=FullScreen,
    }
% You also need to have the packages: cite, citesort, and ifthen
% (ifthen comes with the standard distribution of LaTeX).
% They can be obtained from any CTAN server. Probably where you found
% this.

\newcommand{\degrees}{$^{\circ}$~}

\begin{document}
\DocInput{imac.dtx}
\end{document}
%</driver>
%\fi
%
%%%% Don't want date printed
%\date{}
%
%\title{\fontsize{16pt}{19pt}\selectfont\bfseries IMAC version 2.0: A \LaTeX\ PACKAGE FOR WRITING
%PAPERS IN
%INTERNATIONAL MODAL ANALYSIS CONFERENCE (IMAC) FORMAT}
%
%
%\author{\vspace{.25in}\\
%\textbf{Joseph C. Slater}\\
%}
%\maketitle
%
%%%%%%%%%%%%%%%%%%%%%%%%%%%%%%%%%%%%%%%%%%%%%%%%%%%%%%%%%%
%
%\begin{abstract}
%Just like in any other \LaTeX\  document, you begin the abstract, make
%the title, and define the author as shown in the example. Not that
%IMAC requires that the author names be in bold, so you have to do this
%yourself. In addition, I haven't yet hacked the style definition for
%the title. Maybe I will later, but that's a minor inconvenience.
%
%For feedback and support, please see the \href{https://github.com/josephcslater/IMAC}{IMAC style home page}.
%\end{abstract}
%
%\keywords{Keyword one \textbullet\ Keyword two \textbullet\ Keyword three}
%
%
%\subsubsection*{Nomenclature}
%\begin{tabular}{lll}
%\matr{M} &&  matrix\\
%
%\dmat{M}&& diagonal matrix\\
%
%\vect{\Phi} && vector \\
%
%\elem{M_{11}}&& single element of \matr{M}\\
%
%\pnorm{x} && p-norm of \vect{x}\\
%
%$\begin{bmatrix}\matr{M}_{11}&:&\matr{M}_{12}\\
%	\hdotsfor[2]{1}&:&\hdotsfor[2]{1}\\
%	\matr{M}_{21}&:&\matr{M}_{22}
%\end{bmatrix}$
%&& partitioned matrix\\
%
%\end{tabular}
%
%\section{Changes}
%As of this release, versions are identified by release date
%(\texttt{YYYY-MM-DD}) rather than by an incrementing version number.
%Older entries below retain their original version numbers since
%precise individual dates for each of them could not be reliably
%recovered from the repository history; the closest available
%commit dates are noted.
%
%2018-03-26: Initial public release (version 1.0, later also labeled
%1.0.1/1.0.2), including the LPPL license and hyperlinks to the main
%document.
%
%2021-12-08: (version 1.0.3) Fix minor typos in document.
%
%2026-10-07: Realigned the style with the official IMAC Word Template
%(\emph{Word Template.dotx}, distributed via the IMAC/SEM (Society for
%Experimental Mechanics) author toolkit) so that \LaTeX\ and Word
%submissions now produce matching formatting. Specific changes:
%\begin{itemize}
%\item Switched body and heading fonts from sans-serif (Helvetica/Computer
%Modern Sans) to Times New Roman-compatible fonts (via \texttt{mathptmx}).
%\item Changed from a two-column US Letter layout to a single-column
%A4 layout, with margins matching the Word template (top 0.825in,
%bottom 1in, left/right 0.715in).
%\item Section, subsection, and subsubsection headings are now bold,
%sentence case, instead of bold, all-caps.
%\item The \texttt{Abstract}, \texttt{References}, \texttt{Table}, and
%\texttt{Program} labels are now bold, sentence case, instead of
%bold, all-caps.
%\item In-text citations are now plain, non-superscript, bracketed
%numbers (e.g.\ \texttt{[1, 2, 3]}), matching the Word template, instead
%of bracketed superscripts. Citation ranges are no longer compressed.
%\item Added a \verb*z\keywordsz command to typeset a Keywords line
%after the abstract, matching the Word template.
%\item Paragraph spacing (\texttt{\textbackslash parskip}) reduced to
%8pt to match the Word template's default paragraph spacing.
%\item Removed unused eleven/ten/eight-point \texttt{\textbackslash font}
%definitions, the obsolete \texttt{helvet} package instructions, and
%the duplicate/stray package loading; added \texttt{\textbackslash
%ProvidesPackage} for proper version reporting (now also date-based).
%\item Fixed a hyperlink color inconsistency: citations now use
%\texttt{citecolor=blue}, matching the rest of the document's links
%(previously citations defaulted to hyperref's green).
%\end{itemize}
%
%\section{Typing Your Document}
%\indent First thing to note is that the IMAC style indents the first
%line of each paragraph, including the paragraph immediately following
%a section header. You can override the \LaTeX\ default of not
%indenting by using the \indent command each time after a
%\verb*a\sectiona command or by using the package \verb*zindentfirstz
%which can be downloaded from any CTAN location.
%
%The second thing which you may need to do is obtain and install the
%packages: \verb*zcitez, \verb*zifthenz, \verb*zamsmathz, and \verb*zmathptmxz
%(for Times-compatible text and math fonts, matching the official Word
%template). These are all standard CTAN packages and should already be
%part of any complete \LaTeX\ distribution.
%
%The PDF\textregistered\ file \verb*zimac.pdfz can be viewed for comparison to the results
%you obtain from \LaTeX ing this document.
%
%A few macros have been defined below for conforming to the
%IMAC\cite{ewins} notation convention (See Table \ref{dm}). If you use them, you can simply redefine the
%macros according to the journal requirements when the time comes for submission.
%I've defined only those that I thought were either difficult or do not
%conform to normal textbook standards. You may find it useful to define
%your own macros at the end of the file \verb*zimac.styz, but please
%make sure that you don't delete them the next time you update!
%
%In Table \ref{dm} you'll also see some code for making a partitioned
%matrix. Sorry, but I don't know how to turn this into an
%environment. If you figure it out, let me know and I'll incorporate it.
%The colons form the horizontal delimiters and the \verb*z\hdotsforz
%command forms the vertical delimiters. The first argument represents
%a spacing of the dots, and the second required argument is the number
%of columns that the dotted line should span.
%
%\begin{table*}
%	\begin{tabular}{lll}
%	\matr{M} &\verb*c\matr{M}c&  matrix\\
%
%	\dmat{M}&\verb*c\dmat{M}c& diagonal matrix\\
%
%	\vect{\Phi} & \verb*a\vect{\Phi}a & vector \\
%
%	\elem{M_{11}}&\verb*z\elem{M_{11}}z& single element of \matr{M}\\
%
%	\pnorm{x} &\verb*a\pnorm{x}a & p-norm of \vect{x}\\
%%%%%  Look between these lines for the partitioned matrix code
%	$\begin{bmatrix}\matr{M}_{11}&:&\matr{M}_{12}\\
%		\hdotsfor[2]{1}&:&\hdotsfor[2]{1}\\
%		\matr{M}_{21}&:&\matr{M}_{22}
%	\end{bmatrix}$
%%%%%  Look between these lines for the partitioned matrix code
%%%%%  The stuff below this is for displaying the code. Ignore it.
%	&\begin{minipage}{3.5in}
%	{\verb*z\begin{bmatrix}z\\\verb*z\matr{M}_{11}&:&\matr{M}_{12}\\z\\
%\verb*z\hdotsfor[2]{1}&:&\hdotsfor[2]{1}\\z\\
%\verb*z\matr{M}_{21}&:&\matr{M}_{22}z\\\verb*z\end{bmatrix}z}
%\end{minipage}
%	& partitioned matrix\\
%	\end{tabular}
%	\caption{\label{dm}Defined macros}
%
%\end{table*}
%
%Please contact me at \verb*zjoseph.c.slater@gmail.comz if you find bugs in this.
%I'll do my best to fix them in a timely fashion. Please don't contact me with respect to general \LaTeX\
%questions.
%For help with \LaTeX\ please consult Lamport\cite{lamport}, Goossens
%et al.\cite{goossens}, and/or Kopka and Daly\cite{kopka}. Additional
%resources are available through the \TeX\ newsgroup \verb*zcomp.text.texz,
%the
%CTAN archives at \verb*zhttp://www.ucc.ie/cgi-bin/ctanz, and the \TeX
%users group (TUG) home page (\verb*zhttp://www.tug.org/z).
%
%From here on is just some examples to give you some
%continuing examples of how to do things. The table would have been
%better off if I had used the \verb*zhhlinez package.
%
%
%%%%%%%%%%%%%%%%%%%%%%%%%%%%%%%%%%%%%%%%%%%%%%%%%%%%%%%%%%
%
%\section{Here is a section}
%
%\indent In bladed disk assemblies, the disk acts as a coupling device between the
%blades.  As the stiffness of the disk increases, blade coupling
%decreases. It has been shown that weak interblade coupling leads to
%high levels of mode localization when blades are
%mistuned.
%
%
%\begin{figure}
%\begin{center}
%\fbox{There once was a figure here}
%\end{center}
%\caption{\label{undeformed}Here is a figure.}
%\end{figure}
%
%\subsection{This is a subsection}
%
%\indent All models used in this study were variations of the symmetric,
%constant stiffness system referred to as the baseline model (Figure
%\ref{undeformed}).  Yada yada.
%
%
%
%\begin{table}
%\begin{center}
%\begin{tabular}{|c|c|c|c|}
%\hline
%    & Random 1      &     Random 2  &    Random 3\\
%\cline{2-4}
%\raisebox{2ex}{Blade Number}  & $10^{-3}$ Kg   &    $10^{-3}$ Kg &    $10^{-3}$  Kg\\
%\hline
%\hline
%
%1   &    0.0168  & 0.3343  & 0.3413\\
%2    &   0.0260  & 0.2867  & 0.4431\\
%3    &   0.2579  & 0.4529  & 0.3710\\
%
%\hline
%
%\end{tabular}
%\end{center}
%\caption{\label{random} Simple table.}
%\end{table}
%
%... and that's it! If you can't get it to work, I can be reached at
%\verb*zjoseph.c.slater@gmail.comz.
%
%
%
%
%\section*{Acknowledgments}
%Note that for the Acknowledgments you need to use the
%\verb*z\sectionz command in the starred form to avoid getting the
%section number.
%
%Thanks to Leslie Lamport\cite{lamport}, Goossens, Mittelbach, and
%Samarin\cite{goossens} and all others  who've built \LaTeX\ into what it is
%today.
%
%\section*{Copyright}
%
%Copyright (C) 1998 Joseph C. Slater
%
%\bibliographystyle{imac}
%
%\bibliography{imac}
%
%
%
%\StopEventually{}
%\section{The Code}
%\iffalse
%    \begin{macrocode}
%<*imac.sty>
%    \end{macrocode}
%\fi
% This is imac.sty for producing IMAC-format (International
% Modal Analysis Conference)  articles using LaTeX2e 
% Supported by Joseph C. Slater
% joseph.c.slater@gmail.com
% https://josephcslater.github.io
%
% Copyright (C) 1999 Joseph C. Slater
%
% This program is free software; you can redistribute it and/or
% modify it under the terms of the GNU General Public License
% as published by the Free Software Foundation; either version 2
% of the License, or (at your option) any later version.
%
% This program is distributed in the hope that it will be useful,
% but WITHOUT ANY WARRANTY; without even the implied warranty of
% MERCHANTABILITY or FITNESS FOR A PARTICULAR PURPOSE.  See the
% GNU General Public License for more details.
%
%
% Modified from the ASME style
% written by   Srinivas S. Sripada
%              Research Fellow,
%              Dept. of Mechanical Engg. & Applied Mechanics
%              University of Pennsylvania
%              Philadelphia
%
%%%
%%% Version 2026-10-07
%%% (versioned by release date rather than a version number; see the
%%% "Changes" section of imac.tex for the full history)
%%%

\NeedsTeXFormat{LaTeX2e}
\ProvidesPackage{imac}[2026/10/07 IMAC conference paper format (aligned with the official Word template)]

\RequirePackage{ifthen}
\RequirePackage{amsmath}
\RequirePackage[nocompress]{cite}
\RequirePackage{geometry}

% The official IMAC Word Template.dotx uses Times New Roman body text,
% a single (non-two-column) layout, and A4 paper, so this style now
% mirrors that instead of the old sans-serif/two-column convention.
\RequirePackage{mathptmx} % Times-compatible text and math fonts

%% define the section headings
%% Word template headings are simply bold, sentence/title case (not
%% small caps or all caps) at body text size.
\setcounter{secnumdepth}{1}

\renewcommand{\section}{\@startsection
{section}{0}{0 in}{1\baselineskip}{.001 em}{\normalfont\bfseries}}

\renewcommand{\subsection}{\@startsection
{subsection}{1}{0in}{1\baselineskip}{0.001 em}{\normalfont\bfseries}}

\renewcommand{\subsubsection}{\@startsection
{subsubsection}{0}{0 in}{1\baselineskip}{.001 em}{\normalfont\bfseries}}


%% define headings for Bibliography listing, figures, tables, programs
%% (sentence case, bold, matching the Word template's "Abstract" /
%% "References" labels rather than the previous all-caps sans-serif style)

\renewcommand{\refname}{\noindent {\bfseries References}}
\renewcommand{\abstractname}{\noindent {\bfseries Abstract}}
\renewcommand{\@makecaption}[2]{\vspace{10pt}%
{\begin{center}{\bfseries #1: #2}\end{center}}%
}

\def\tablename{\bfseries Table}
\def\Programname{\bfseries Program}

% A Keywords line, as used after the abstract in the Word template, e.g.:
%   \keywords{Keyword one \textbullet\ Keyword two \textbullet\ Keyword three}
\newcommand{\keywords}[1]{\par\noindent{\bfseries Keyword} \quad #1\par}

% Page geometry matching the Word template (A4 paper; see Word Template.dotx
% word/document.xml sectPr: pgSz 11907x15876 twips, pgMar top=1188,
% bottom=1440, left=right=1030 twips).
\geometry{a4paper, top=0.825in, bottom=1in, left=0.715in, right=0.715in}

\setlength{\itemsep}{0pt}
\setlength{\parskip}{8pt} % matches Word docDefaults spacing w:after="160" twips
\setlength{\parindent}{0in}
\renewcommand{\textfraction}{0.2}
\renewcommand{\floatsep}{18pt}
\renewcommand{\textfloatsep}{18pt}
\renewcommand{\intextsep}{.1875 in}
\renewcommand{\floatpagefraction}{.7}

% The Word template's in-text citations are plain inline brackets
% (e.g. "[1, 2, 3]"), not bracketed superscripts, and ranges are not
% compressed. The `cite` package (loaded with [nocompress] above)
% already gives non-superscript bracketed citations; \citepunct just
% adds the space after the comma to match the Word template's style.
\renewcommand{\citepunct}{,\ }

% We don't want any headers or footers. IMAC will put on what they like.
\pagestyle{empty}

% The following prevents the maketitle from putting a page number on 
% page 1
\let\ps@plain=\ps@empty

\newcommand{\vect}[1]{\ensuremath{\left\{#1\right\}}}
\newcommand{\matr}[1]{\ensuremath{\left[#1\right]}}
\newcommand{\elem}[1]{\ensuremath{\left(#1\right)}}
\newcommand{\dmat}[1]{\ensuremath{[\hbox{\textsf{\`}}#1\raisebox{-5pt}{\hbox{\textsf{\`}}}]}}
\newcommand{\pnorm}[1]{\ensuremath{\lVert#1\rVert_{p}}}

\endinput



%\iffalse
%    \begin{macrocode}
%</imac.sty>
%    \end{macrocode}
%\fi
%\Finale
\endinput
